\chapter{Zehntausend Saisons: Simulation und Monte Carlo}\label{kap:simulation}

\section{Eine simulierte Saison}

Eine Iteration spielt die ganze Saison durch
(\datei{league-simulator-rust/src/simulation/season.rs:10--62}). Sie beginnt
jedes Mal bei den ELO-Werten zum Saisonstart und geht den Spielplan in
Listenreihenfolge durch. Für jedes Spiel gibt es drei Fälle:

\begin{enumerate}
  \item \textbf{Das Ergebnis fehlt.} Die Tore werden mit dem Tormodell
        gezogen und eingetragen. \emph{Danach wird die ELO aktualisiert}, als
        wäre das Spiel wirklich so ausgegangen.
  \item \textbf{Das Ergebnis liegt vor.} Es zählt für die Tabelle, und die ELO
        wird mit dem echten Ergebnis aktualisiert.
  \item \textbf{Das Ergebnis liegt vor, das Spiel ist gewertet} (Kennzeichen
        \code{elo\_neutral}). Es zählt für die Tabelle, die ELO bleibt
        unverändert.
\end{enumerate}

Die ELO ist innerhalb einer simulierten Saison also \emph{dynamisch}. Eine
Mannschaft, die in einem Durchlauf eine Siegesserie hinlegt, geht stärker in
die folgenden simulierten Spiele. Das bildet ab, dass sich auch die echte
Stärkeschätzung im Lauf der Saison bewegen wird, und verbreitert die
Verteilung der Endplätze gegenüber einem Modell mit festen Stärken.

\subsection*{Reihenfolge der Spiele}

Die R-Seite sortiert den Spielplan chronologisch, bevor er an Rust geht
(\datei{RCode/transform_data.R:271--291}): nach Anstoßzeit, bei Gleichstand
in der Reihenfolge der API, Spiele ohne Termin ans Ende. Gespielte Spiele
werden so in jeder Iteration in derselben Reihenfolge nachgespielt und führen
zu denselben ELO-Werten. Die simulierten Spiele folgen danach in der Reihenfolge
des Spielplans.

\section{Statusgruppen}

Die Daten kommen von API-Football. Jedes Spiel trägt einen Statuscode, und
das Modell fasst die Codes in Gruppen zusammen
(\datei{RCode/league_details.R:56--68},
\datei{RCode/transform_data.R:357--374}):

\begin{center}
\small
\begin{tabularx}{\textwidth}{@{}llXcc@{}}
\toprule
Gruppe & Codes & In der Simulation & ELO-Walk & Tabelle \\
\midrule
beendet    & FT, AET, PEN & Ergebnis fest & ja & ja \\
gewertet   & AWD, WO      & Ergebnis fest, \code{elo\_neutral} & nein & ja \\
live       & 1H, HT, 2H, ET, BT, P, SUSP, INT, LIVE & wird simuliert & nein & nein \\
verschoben & PST, CANC, TBD, ABD & wird simuliert & nein & nein \\
offen      & alle übrigen, vor allem NS & wird simuliert & nein & nein \\
\bottomrule
\end{tabularx}
\end{center}

Laufende Spiele fließen erst mit dem Abpfiff ein. Der Zwischenstand bleibt
unberücksichtigt, und die Website sagt das in der Live-Sektion ausdrücklich.
Simuliert werden nur Hauptrundenspiele. Runden mit Bezeichnungen wie
\enquote{Relegation}, \enquote{Play-off} oder \enquote{Final} filtert
\datei{RCode/round_filter.R} heraus.

\section{Die Abschlusstabelle und ihre Reihenfolge}\label{sec:tiebreak}

Am Ende jeder Iteration entsteht eine Tabelle mit drei Punkten für einen
Sieg, einem für ein Remis und keinem für eine Niederlage. Sortiert wird
stabil (\datei{season.rs:92--168}):

\begin{enumerate}
  \item Punkte, absteigend,
  \item Tordifferenz, absteigend,
  \item erzielte Tore, absteigend,
  \item bei völligem Gleichstand die Eingabereihenfolge der Mannschaften.
\end{enumerate}

Die Eingabereihenfolge ist alphabetisch nach Kürzel
(\datei{transform_data.R:383--387}). Bei völligem Gleichstand gewinnt also
systematisch das alphabetisch frühere Kürzel. Einen direkten Vergleich wertet
das Modell nicht aus. Bei zehntausend Durchläufen und ganzzahligen Toren ist
ein Gleichstand in allen drei Kriterien selten, kommt aber vor.

Die Engine kann Startwerte für Punkte und Tore übernehmen
(\code{adj\_points}, \code{adj\_goals} und weitere). Im Hauptlauf sind sie
null. Reale Punktabzüge, etwa nach Lizenzverstößen, werden nicht modelliert,
weil es dafür keine Datenquelle gibt. Genutzt werden die Startwerte nur für
den Malus-Lauf der Zweitvertretungen (Abschnitt~\ref{sec:malus}).

\section{Monte Carlo}

Der Rust-Server wiederholt die Saison $N = \num{10000}$-mal
(\datei{league-simulator-rust/src/monte_carlo/mod.rs}). Es sei $C_{t,p}$ die
Zahl der Iterationen, in denen Mannschaft $t$ auf Platz $p$ landet. Die
angezeigte Wahrscheinlichkeit ist der Anteil:
\begin{equation}\label{eq:mc}
  \boxed{\;\hat P(t \text{ auf Platz } p) = \frac{C_{t,p}}{N}\;}
\end{equation}
Jede Zeile und jede Spalte dieser Matrix summiert sich zu 1. Die Zeilen
sortiert der Server nach dem erwarteten Platz
$\bar p_t = \sum_p p\,\hat P_{t,p}$, aufsteigend. Diese Matrix ist die
Heatmap auf jeder Ligaseite.

\begin{description}[font=\sffamily\bfseries\small]
  \item[Zufallszahlen] Jede Iteration erhält einen eigenen 64-Bit-Startwert
        aus dem Thread-Zufallsgenerator des Betriebssystems und erzeugt daraus
        einen eigenen Generator (\code{StdRng::seed\_from\_u64}). Zwei Läufe
        mit denselben Daten liefern deshalb leicht unterschiedliche Matrizen.
        Eine deterministische Variante mit festem Startwert gibt es nur für
        Tests.
  \item[Parallelität] Die Iterationen laufen über \code{rayon} parallel
        (\code{par\_iter().fold().reduce()}). Jeder Thread zählt in eigene
        Zähler, die am Ende addiert werden. Ganzzahlige Addition ist
        kommutativ, die Reihenfolge der Threads ändert das Ergebnis also nicht.
  \item[Laufzeit] Rund 42~ms für \num{10000} Iterationen einer Liga mit 18
        Mannschaften und vollem Restspielplan, gemessen am 7.~September 2026.
        Alle zehn Ligen zusammen brauchen etwa eine halbe Sekunde.
\end{description}

\subsection*{Wie genau sind die Zahlen?}

Der Standardfehler eines geschätzten Anteils ist
\begin{equation}
  \operatorname{SE}(\hat P) = \sqrt{\frac{\hat P (1 - \hat P)}{N}} \le
  \sqrt{\frac{0{,}25}{\num{10000}}} = 0{,}5\ \text{Prozentpunkte}.
\end{equation}
Die Methodik-Seite fasst das allgemeinverständlich zusammen: 35 gegen 37
Prozent heißt gleichauf. Ein Häkchen (\checkmark) bedeutet, dass in keiner
der zehntausend Simulationen etwas anderes herauskam. Das ist etwas anderes
als \enquote{rechnerisch sicher}.

\section{Von der Matrix zu den Spalten der Website}

Die Spalten neben der Heatmap, etwa \enquote{Meister},
\enquote{Champions League} oder \enquote{Abstieg}, summieren Plätze
(\code{groupResultsDF}, \datei{RCode/render_helpers.R:49--77}):
\begin{equation}
  P(t \in \text{Gruppe } [l, u]) = \sum_{p=l}^{u} \hat P_{t,p}.
\end{equation}
Negative Grenzen zählen von unten: $-1$ ist der letzte Platz. Sie werden erst
zur Renderzeit gegen die Ligagröße aufgelöst, weil Ligen ihre Größe ändern
können (die Frauen-Bundesliga wuchs 2025 von 12 auf 14 Mannschaften).
Tabelle~\ref{tab:panels} in Kapitel~\ref{kap:ligaseite} listet alle Gruppen
je Liga.

\subsection*{Zusatzzählung für die 3.~Liga}

Für die Regionalligen ist nicht nur wichtig, wie wahrscheinlich ein
Drittligist absteigt, sondern auch \emph{wohin}. Er fällt in die Regionalliga
seiner Region zurück. Deshalb zählt der Server für die 3.~Liga zusätzlich,
wie viele der vier Absteiger aus jeder Staffelregion kommen. Es sei $G_{g,j}$
die Zahl der Iterationen mit genau $j$ Absteigern aus Region $g$
(\datei{monte_carlo/mod.rs:133--147}). Auf diesen Zählungen baut die
Abstiegskopplung auf (Abschnitt~\ref{sec:kopplung}).
